\section{Experiments}
\label{sec:experiments}

\subsection{Data and evaluation protocol}
\label{sec:exp-data}

The main experiments use the mobile BCI dataset of \citet{lee2021mobile}: 24 participants recorded with 32 scalp and 14 around-the-ear EEG electrodes plus four periocular electrodes while standing, walking, or running on a treadmill and performing two brain--computer interface paradigms (an ERP oddball and an SSVEP task). We represent each recording as 46 EEG channels and two bipolar EOG derivations (vertical and horizontal) at 100 Hz, with 5.12-s model inputs (512 samples), and use three ambulatory speed sessions per participant with both paradigms, giving up to six session--task cells each. The ambulatory setting is deliberate: it is the regime where ocular and motion contamination is largest and where a deployable artifact-removal system matters most. Within each cell, a continuous 120-s prefix supplies the calibration and all evaluation windows occur later in the record. Fifteen participants form the development cohort, split into five participant-held-out folds of nine training, three validation, and three test participants. Eight further participants form the held-out cohort: they were excluded from every modeling and analysis decision, and were evaluated exactly once, after the fold models, the population registry, and the comparison had been frozen, with no iteration on the outcome. The source study was approved by the Institutional Review Board of Korea University (KUIRB-2019-0194-01); the recordings are public under a CC0 license, and this work is a secondary analysis of the de-identified data \cite{lee2021mobile}.

Reconstruction accuracy is measured on paired episodes built entirely from real signals. A low-ocular-activity EEG segment serves as the reference $\mathbf{x}$, a separately sampled real ocular waveform is passed through the recipient's later-record propagation matrix (estimated from 150--270 s, never an input to any model) to form the contamination, and the sum is the observation. Reference, ocular waveform, and recipient calibration come from three different participants, so the construction isolates the propagation pathway while keeping real EEG and real EOG throughout. Injected amplitudes use base gains of 0.35, 0.70, or 1.15 with a uniform factor from 0.85 to 1.15, and 15\% of episodes contain no injected artifact; eight episodes are evaluated per participant across the available session--task cells. The paired endpoint is the relative root-mean-square error, $\rrmse = \lVert \widehat{\mathbf{x}} - \mathbf{x} \rVert_2 / \lVert \mathbf{x} \rVert_2$, computed per episode and averaged within participant; improvements are reported as $\rrmse$ differences in which positive values favor the first-named method. Natural recordings have no clean target and are analyzed separately on four windows per cell placed from 300 s onward, with three endpoints: EOG attenuation in decibels, the reduction of EOG--EEG coherence, and low-EOG retention, the proportion of the recorded waveform preserved in segments with little EOG activity. Paired and natural outcomes are never combined into one score.

A supporting experiment uses EEGEyeNet \cite{kastrati2021eegeyenet} to test whether the benefit of matched calibration depends on its matrix interface or on the information it conveys. A compact deterministic network receives an injected 46-channel episode and one of three calibration representations: participant-matched, zero, or mismatched. For each training-pool size from 30 to 259 participants the network is trained for 20{,}000 AdamW updates, then frozen, and a 32-dimensional representation is fitted per evaluation participant against paired clean calibration targets and tested on disjoint query episodes. Because the representation is fitted with clean targets, this experiment upper-bounds what calibration information can deliver through a learned interface; it is not part of the main method.

\subsection{Comparison conditions}
\label{sec:exp-comparisons}

All diffusion variants share one population network and identical evaluation instances; they differ only in what the guide and calibration features contain (Table~\ref{tab:variants}). \emph{Matched calibration} is the proposed mode: the guide is the evaluation EOG propagated through the user's own shrunk matrix. \emph{Population calibration} replaces the matrix by the recipient-excluded pooled estimate, testing whether an individual estimate is needed at all. \emph{Mismatched calibration} uses another participant's matrix from the same session--task stratum, testing specificity. The \emph{unguided control} withholds the waveform guide while keeping the calibration features, isolating the contribution of the explicit guide; it is also the operating mode the fallback of Section~\ref{sec:method-calibration} reverts to. \emph{Shuffled EOG} destroys the temporal alignment between EOG and EEG. The ablation matrix of Section~\ref{sec:res-paired} additionally crosses the guide factor with the calibration-feature factor (population features under a matched guide and vice versa) and adds an \emph{unshrunk} arm whose guide uses the raw 120-s ridge estimate ($\lambda=1$), quantifying what shrinkage buys. Outside the diffusion family, we evaluate raw EEG, linear EOG regression from the same calibration \cite{gratton1983ocular}, a one-step deterministic network with the same inputs, an on-line variant in which the propagation matrix is tracked by recursive least squares during evaluation instead of being calibrated beforehand, and, for interval quality, an ensemble of the deterministic network. Three literature anchors run on the same protocol as context rows rather than contests: ICA with EOG-correlation component removal \cite{jung2000bss,pion2019iclabel} fitted on each continuous record (the 5.12-s paired episodes are too short for a stable 46-channel decomposition, a constraint we report as such), ASR calibrated on the same 120-s prefix \cite{mullen2015asr,chang2020asr}, and a calibrated eye-subspace projection in the style of SGEYESUB \cite{kobler2020sgeyesub} that projects out the rank-2 column space of the same calibration ridge fit.

\begin{table*}[t]
\caption{Comparison conditions and the information each receives. All diffusion modes use the same trained population network.}
\label{tab:variants}
\centering
\small
\begin{tabularx}{\textwidth}{@{}p{.19\textwidth}Yp{.20\textwidth}p{.24\textwidth}@{}}
\toprule
Condition & Propagation source & Guide / features & Question answered \\
\midrule
Matched calibration & Own calibration segment & Propagated EOG; own features & Proposed method \\
Population calibration & Recipient-excluded pool & Propagated EOG; population features & Is the individual estimate needed? \\
Mismatched calibration & Another participant & Propagated EOG; donor features & Is the estimate specific to its owner? \\
Unguided control & Own calibration segment & No guide; own features & What does the explicit guide add? \\
Shuffled EOG & Own calibration segment & Misaligned guide; own features & Does timing matter? \\
Linear EOG regression & Own calibration segment & Direct subtraction & Classical reference \\
Deterministic network & Own calibration segment & One-step reconstruction & Learned point-estimate reference \\
On-line adaptation & Recursive least squares during use & Propagated EOG & Can calibration be replaced on-line? \\
Deterministic ensemble & Own calibration segment & Ensemble spread & Interval-quality reference \\
\bottomrule
\end{tabularx}
\end{table*}

\subsection{Adaptation-loop and downstream protocols}
\label{sec:exp-loop}

Three further protocols characterize the system as a loop. \emph{Calibration duration}: the calibration prefix is truncated to 30, 60, 90, or 120 s and scaling, ridge, and shrinkage are recomputed closed-form; each duration is read twice, once with the reliability rule active (the system reading; 30 s is rejected by construction since the four sub-block statistics require 60 s) and once with the rejection floor bypassed (the raw shrinkage curve). \emph{Operator lifetime}: the propagation matrix is re-estimated closed-form on successive 120-s windows across each record and its relative displacement from the calibration estimate is tracked; the natural-recording endpoints are additionally stratified by window position, and fresh paired episodes drawn from early, mid, and late record thirds (same construction, restricted window starts) read the paired gain against elapsed time. \emph{Downstream utility}: the dataset's own tasks are decoded from each method's output with deliberately shallow, deterministic decoders --- the three-class SSVEP trials (5 s, one model window) by canonical correlation against reference sinusoids and their second harmonics at the three stimulus frequencies, frozen from a probe cell before any evaluation, over thirteen occipito-parietal channels; and the ERP oddball ($-0.2$ to $0.8$ s epochs, standard:deviant 4:1) by shrinkage-LDA under within-participant five-fold cross-validation with AUC as the endpoint, plus an ERP-preservation check correlating each method's low-contamination trial average with the raw one. Trials beginning inside the calibration prefix are excluded, and SSVEP contrasts are additionally stratified by within-participant tertiles of trial EOG energy, with the interpretation grid for all downstream outcomes frozen before compute.

\subsection{Implementation and statistical analysis}
\label{sec:exp-implementation}

The shared network is a four-scale one-dimensional U-Net (channel widths 32--128, four-head self-attention at the bottleneck, 128-dimensional calibration context, 965{,}085 parameters) predicting clean EEG over 1{,}000 linear diffusion levels ($\beta_1=10^{-4}$, $\beta_{1000}=0.02$). Each fold and seed trains for 80{,}000 AdamW updates (batch 8, learning rate $10^{-4}$, weight decay $10^{-4}$, gradient clipping 1.0, exponential moving average 0.999). During training the guide is dropped with probability 0.30 and the calibration features are replaced by their population values with probability 0.20, so the unguided and fallback modes are in-distribution at evaluation time. Development point estimates use five folds, three seeds, and 50 deterministic DDIM steps with initial noise shared across conditions. The held-out evaluation uses one frozen model per fold, a population registry built from the 15 development participants, and the mean of the five fold predictions with a common noise draw per episode. The interval analysis uses $K=32$ stochastic trajectories per development fold--seed cell with the temperature chosen without the evaluated fold. The paired comparisons reported in Section~\ref{sec:results} were run with an earlier version of the rejection rule that substituted population calibration instead of withholding the guide; the two rules act identically on every cell whose calibration is accepted, and rerunning the development comparison under the final rule left the headline effect unchanged (Section~\ref{sec:res-paired}).

All contrasts are aggregated within participant first: windows, protocols, seeds, and conditions are averaged per participant, and we report participant-level means, medians, counts of positive effects, and 95\% percentile bootstrap intervals from 5{,}000 participant resamples. Interval coverage is computed over channel--time elements within each fold--seed cell and averaged over the 15 cells; the continuous ranked probability score and risk--coverage area are aggregated over episodes. Code, configurations, and stored evaluator outputs will be released on publication.
